\section{Related Work}
\label{sec:related_work}

\paragraph{Dutch Electricity Grid Load Forecasting} 
Recent studies in Dutch electricity grid load forecasting \cite{van_de_sande_enhancing_2024,ashtar_hybrid_2025,curiel_integrating_2025} are comparable to this study; they all seek to predict the national hourly mean system load (in megawatts) for the Netherlands, also sourcing their data from the European Network of Transmission System Operators for Electricity (ENTSO-E).

\cite{van_de_sande_enhancing_2024} investigated the predictability of wintertime energy demand using a dataset spanning from 2009 to 2019. Testing a classical Multiple Linear Regression (MLR) baseline against an ensemble Prophet-LSTM architecture, they found that the Prophet-LSTM model, where Prophet captures trend, seasonality and holiday effects while the LSTM models non-linear residual dependencies, performed best across all evaluated metrics: RMSE, MAPE and PCC. The ensemble attained its highest performance ($r=0.93$), enabling accurate forecasts up to 180 days in advance, when trained purely on endogenous features. When evaluating exogenous variables iteratively, the authors fetched local weather data from the KNMI, climate indices (such as the North Atlantic Oscillation, Sea Surface Temperature, and Sea Ice Concentration) from U.S. National Oceanic and Atmospheric Administration (NOAA) and the Copernicus Climate Data Store, and Dutch energy prices and renewable energy ratios from CBS. They found that while adding CBS energy price predictors yielded modest improvements, incorporating climate change-related predictors collectively reduced model performance due to severe multicollinearity.

Extending the analysis timeframe into periods characterized by high market volatility, \cite{ashtar_hybrid_2025} utilized a dataset ranging from 2009 to 2023, thus explicitly encompassing both structural breaks introduced by the COVID-19 pandemic and the subsequent European energy crisis. The authors compared a classical SARIMAX baseline against a hybrid SARIMAX-LSTM and a Sequence-to-Sequence (Seq2Seq) architecture. To manage the extreme load anomalies introduced by the crisis events, the authors applied Winsorization to cap extreme outliers at the 1st and 99th percentiles. Furthermore, they provided structural context to the models by including a COVID-19 lockdown dummy flag and an energy import ratio, alongside actual consumer energy and gas prices from CBS. Evaluated on metrics including RMSE, MAPE, MAE, nRMSE, and nMAE, the Seq2Seq model trained strictly with engineered temporal features achieved the highest overall accuracy, demonstrating a superior MAPE of 1.88\% and an nRMSE of 2.86\% at a 180-day forecasting horizon. Their study additionally evaluated exogenous predictors—incorporating KNMI weather data, CBS socioeconomic indicators (GDP, population), and CBS renewable energy production data. The study revealed that while integrating renewable energy data substantially improved short-term forecast accuracy, it contributed little to long-term predictions, and socioeconomic indicators exhibited limited overall predictive value.

Finally, \cite{curiel_integrating_2025} bridged the gap between macroeconomic shocks and quantum-enhanced architectures using an extensive dataset spanning 2010 to 2024, capturing the full span of the recent energy crisis. The study evaluated standard Prophet and Prophet-LSTM baselines against a Prophet-QLSTM (a quantum-enhanced architecture), extending the analysis to include XGBoost-based stacked generalizations of the hybrid models. They retained the anomalous crisis observations and relied dynamically on Prophet's inherent changepoint detection and the XGBoost meta-learner to adapt. They employed Bayesian optimization to select features from a pool including KNMI local weather, Copernicus/ERA5 and NOAA climate indices, CBS macroeconomic indicators (including the Consumer Price Index, housing prices, and labor participation), and global commodity futures from the World Bank's Pink Sheet. Evaluated via rolling-origin cross-validation, their results showed that the XGBoost-stacked variants substantially and consistently outperformed non-stacked counterparts across RMSE, MAPE, and PCC metrics. The stacked quantum-inspired QLSTM model achieved statistical parity with the classical stacked LSTM in terms of RMSE and MAPE, despite utilizing approximately 29 times fewer recurrent parameters, though the classical stack retained a marginal edge in PCC. Their feature selection process confirmed that while short-run load variations are dominated by weather effects, the inclusion of selected World Bank commodity futures and CBS economic indicators was useful for stabilizing long-horizon baselines during structural pattern breaks.

While temporal features consistently form the robust core of the best models in all above studies, the reported utility of exogenous variables varies significantly depending on the forecasting horizon and dataset timeframe. Despite instances where external variables introduced detrimental multicollinearity \cite{van_de_sande_enhancing_2024}, specific exogenous features have proven instrumental in navigating severe structural breaks. For example, \cite{curiel_integrating_2025} demonstrated that macroeconomic indicators and global commodity futures were essential for stabilizing long-horizon baselines during the European energy crisis, while \cite{ashtar_hybrid_2025} effectively utilized energy prices, import ratios, and lockdown flags to the same end. Given these mixed findings, our study will investigate other approaches for introducing and evaluating exogenous predictors, including dimensionality reduction and permutation feature importance analysis. 

The methodologies established in these works \cite{curiel_integrating_2025,van_de_sande_enhancing_2024,ashtar_hybrid_2025} demonstrate the necessity of Rolling-Origin Cross-Validation (ROCV) for robust time-series evaluation that respects temporal integrity, and they establish a standard suite of evaluation metrics (e.g., RMSE, MAPE, and PCC). These studies equip our research with a pool of reference baselines optimized specifically for the Dutch grid, allowing us to benchmark our approach against highly competitive architectures such as Sequence-to-Sequence (Seq2Seq) networks \cite{ashtar_hybrid_2025} and Prophet-LSTM ensembles \cite{van_de_sande_enhancing_2024}.

\paragraph{Incorporation of NTL} All the aforementioned approaches from the literature rely on historical load data combined solely with other time-series data, including meteorological data and traditional economic indicators. However, the broader forecasting domain provides compelling examples of incorporating \textbf{multi-modal} data: Nighttime Light (NTL) satellite imagery for electricity grid load forecasting \cite{borunda_machine_2025,r_forecasting_2025,chen_monthly_2025}, and satellite-image fusion in other problem domains \cite{schubnel_solarcrossformer_2025,wang_causal-aware_2025}. NTL imagery has proven to be a highly effective proxy for true electricity demand across numerous geographies. 

In contrast to the research established in other domains \cite{schubnel_solarcrossformer_2025,wang_causal-aware_2025}, the current system grid load prediction/energy demand forecasting research does not incorporate true 2D spatial image processing, such as applying 2D Convolutional Neural Networks (CNNs) or Vision Transformers directly to raw satellite pixel grids. Instead, it employs aggregates such as the total "Sum of Lights" (SOL) or the sum of pixel intensities within predefined administrative boundaries. These derived scalars are then treated as standard 1D time-series covariates \cite{borunda_machine_2025,r_forecasting_2025,chen_monthly_2025}.

Despite bypassing direct spatial computer vision techniques, these aggregate-based methods work well. In \cite{chen_monthly_2025}, the authors summed the NTL pixel values per city. This scalar sum, combined with monthly average temperature data, was fed into a parameterized polynomial regression model to predict urban electric power consumption for the Yangtze River Delta. Their model achieved an out-of-sample Mean Absolute Relative Error (MARE) of 10.38\% for monthly predictions and successfully generated high-resolution spatial distribution maps by applying these scalar relationships to individual grids. 

Similarly, \cite{r_forecasting_2025} extracted a state-wise SOL index to incorporate as an auxiliary temporal feature in deep learning models in India. Although they utilized architectures named CNN-LSTM and ConvLSTM, these models applied 1D temporal convolutions strictly to the aggregated time-series vectors (load and SOL) rather than 2D spatial convolutions on the imagery. Their multivariate Bidirectional LSTM achieved the lowest overall Root Mean Square Error (RMSE) for the country, while their state-wise approach, utilizing k-means state clustering to capture macroscopic spatial variations, yielded MAPE values below 10\% across all regional clusters. 

Finally, \cite{borunda_machine_2025} proposed a machine learning framework to accurately forecast regional electricity consumption one year ahead in Mexico. They operated by forecasting NTL pixel-by-pixel using independent 1D Random Forest regressors, then summing these pixel predictions into municipal aggregates and mapping them to historical load via a logarithmic linear regression. This aggregate approach achieved a high correlation ($R^2 > 0.85$) with official grid load records, and feature importance analysis revealed that the aggregated NTL sum was the overwhelmingly dominant predictive parameter, providing far more predictive power than auxiliary variables such as temperature and GDP per capita \cite{borunda_machine_2025}. Ultimately, while the reduction of NTL imagery to aggregate sums provides a highly reliable and computationally efficient proxy that successfully drives accurate predictions, the native integration of 2D spatial image processing remains underexplored in this problem domain.

\paragraph{Multi-modal attention} Although previous studies \cite{borunda_machine_2025,r_forecasting_2025,chen_monthly_2025} incorporate NTL aggregate data, attention-based mechanisms are not considered. Attention-based models have proven effective in capturing complex inter-modal relationships in other problem domains \cite{schubnel_solarcrossformer_2025,wang_causal-aware_2025}. Research has explored the integration of diverse data sources, such as text and time-series for supply chain demand forecasts \cite{wang_causal-aware_2025} or satellite imagery for solar irradiance forecasts \cite{schubnel_solarcrossformer_2025}. \cite{wang_causal-aware_2025} proposed a Causal-Aware Multimodal Transformer framework that systematically integrates text, time series, and satellite imagery, reducing RMSE by 12.3\% and MAPE by 15.7\% compared to state-of-the-art baselines for a supply chain demand prediction problem. Their causal analysis demonstrated that satellite-derived economic indicators strongly drive long-term forecasting accuracy. \cite{schubnel_solarcrossformer_2025} introduced SolarCrossFormer, a deep learning model that leverages novel graph neural networks combined with cross-attention to fuse satellite imagery with ground-based meteorological time series. The model achieved a normalized mean absolute error of 6.1\% for day-ahead solar irradiance forecasting and exhibited strong operational robustness by producing accurate probabilistic forecasts for unseen locations lacking historical input data via their geographic coordinates. These attention-based mechanisms have demonstrated significant promise in capturing inter-modal relationships and could be adapted to the problem of grid load forecasting.

\paragraph{Research Gap} A gap exists in the current literature: although models have been designed to simultaneously integrate aggregate information from VIIRS NTL imagery, economic indicators, and weather data via other mechanisms, as in \cite{chen_monthly_2025,borunda_machine_2025,r_forecasting_2025}, multi-modal attention-based mechanisms involving the native integration of 2D images have not been adapted for this use case, nor tested for the Dutch context. Furthermore, within the Dutch context, no inclusion/exclusion study has quantified the usefulness of NTL aggregate data alone or time series self-attention for this problem statement. This research aims to bridge these gaps by adapting multi-modal architectures similar to \cite{wang_causal-aware_2025} and \cite{schubnel_solarcrossformer_2025} for the problem at hand, while benchmarking results against established Dutch baselines \cite{van_de_sande_enhancing_2024, ashtar_hybrid_2025} using comparative metrics (RMSE, MAPE, MAE, MASE, PCC) and inclusion/exclusion studies as in \cite{wang_causal-aware_2025} to quantify the predictive gain of the self-attention, NTL 1D aggregate, and 2D NTL/cross-attention components.
